\documentclass[11pt]{report}

\usepackage[pdftex]{geometry}
\geometry{a4paper,left=2.8cm,right=2.8cm,top=1.5cm,bottom=1.5cm,twoside}

\usepackage{amssymb}
\usepackage{enumerate}
\usepackage{graphicx}
	\graphicspath{./Comments/}
\usepackage[many]{tcolorbox}
\usepackage{natbib}
\usepackage{hyperref}
\usepackage{subfigure}
\usepackage{bm,bbm}
\usepackage{float}
\usepackage{siunitx}

\usepackage[utf8]{inputenc}
\usepackage[T1]{fontenc}
\usepackage{xcolor}
\newtcolorbox{reviewbox}[1]{
colback = white!5!white,
colframe = purple!75!black,
fonttitle=\bfseries,
title=#1
}

\newtcolorbox{responsebox}[1]{
	colback = white!5!white,
	colframe = white!75!black,
	fonttitle=\bfseries,
	title=#1
}
\newcommand{\assigned}[1]{\textcolor{blue}{\textbf{#1}}}
%%%%%%%%%%%%%%%%%%%%%%%%%%%%
\begin{document}


\begin{center}
\large{\textbf{Response Letter to Reviewer \# 2}}

\vglue 0.3cm

\huge{ Adaptive Model-Free Tsallis Entropy Estimation for Detecting Non-Fully-Developed Speckle\\ (TGRS-2025-10064)}
\end{center}

%\authors
\begin{center}
\textbf{Janeth Alpala,  Alejandro C.\ Frery, Paolo Gamba, Abraão D. C. Nascimento, Andrea A. Rey }
\end{center}

\date{\today}

%%%%%%%%%%%%%%%%%%%%%%%%%%%%

\vspace{0.5cm}
\noindent Dear Editor
\bigskip

\noindent We are grateful to the reviewers for the valuable comments and the time spent on checking our manuscript. 
We have addressed their comments towards improving the paper contents and presentation. 
Below we detail the changes and updates incorporated into the manuscript in this round of review.

\medskip



\vspace{0.5cm}

%-------------------------------------
\begin{reviewbox}{Comment}
The manuscript proposes an unsupervised hypothesis-testing framework for detecting departures from the fully developed speckle hypothesis in SAR intensity imagery. The core contribution combines a nonparametric estimator of Tsallis entropy (with bootstrap bias correction) and an adaptive windowing scheme driven by the local coefficient of variation.

The topic fits TGRS, and the technical derivations appear sound. However, I have key concerns about (i) the novelty relative to the authors' recent related publications and (ii) whether the current validation isolates the benefit of using Tsallis entropy, as opposed to the benefit of adaptive windowing. Specific comments are provided below.
\end{reviewbox}

\begin{responsebox}{Response}

\assigned{Revision assigned to Janeth}

\end{responsebox}

\vspace{1em}

%-------------------------------------
\begin{reviewbox}{Comment 1}
Concerns Regarding Novelty and Incremental Contribution

The manuscript is closely related to the authors' recent works, specifically the Shannon-entropy test ([18], [30]) and the Rényi-entropy test ([19]). In this submission, the main changes appear to be: (a) replacing Shannon/Rényi with Tsallis entropy and (b) adding an adaptive window selection step.

- Relationship to prior work

Moving from Shannon $\rightarrow$ Rényi $\rightarrow$ Tsallis is mathematically reasonable, but it may read as a straightforward extension unless the paper clearly demonstrates a distinct statistical or practical advantage of Tsallis entropy in this specific SAR setting. At the moment, the paper appears to follow a very similar testing pipeline to [18] and [19], with a different entropy definition.

- Recommendations

Please clearly state—and quantitatively demonstrate—why Tsallis entropy is preferable here. In particular, does the extra parameter $\lambda$ lead to a meaningful gain (e.g., better separation of texture levels, improved detection in certain heterogeneity regimes, or stronger robustness) compared with Rényi and Shannon approaches?
\end{reviewbox}

\begin{responsebox}{Response}

\assigned{Revision assigned to Janeth}

\end{responsebox}

\vspace{1em}

%-------------------------------------
\begin{reviewbox}{Comment 2}
Missing benchmarking against entropy-based baselines

In Section IV, the experiments mainly compare adaptive windows versus fixed windows. This is not sufficient to support the central claim about the Tsallis-based statistic.

- Missing baselines

Since Shannon- and Rényi-based tests already exist in [18] and [19], the paper should compare the proposed Tsallis test against those baselines directly.

- Isolating the source of improvement

Without a comparison such as ``Adaptive Tsallis vs.\ Adaptive Rényi vs.\ Adaptive Shannon'', it is not possible to tell whether improvements come from the Tsallis statistic or simply from adaptive windowing (which could be applied to any entropy measure).

- Recommendations

At minimum, please provide a threshold sweep (e.g., $p$-threshold from $0.01$ to $0.20$) for Tsallis, Rényi, and Shannon under the same windowing protocol and on the same ROIs (ideally for both fixed and adaptive windows). If space permits, ROC curves and/or precision--recall curves would provide a complementary, threshold-free view.
\end{reviewbox}

\begin{responsebox}{Response}

\assigned{Revision assigned to Janeth}

\end{responsebox}

\vspace{1em}

%-------------------------------------
\begin{reviewbox}{Comment 3}
Calibration with adaptive window selection

I am concerned about whether the reported $p$-values remain well calibrated when the window size is chosen adaptively (see the adaptive-windowing part around Eq.~(15)).

- Data-dependent window selection

The manuscript reports null behavior (Table I) based on Monte Carlo simulations for fixed sample sizes. In the proposed pipeline, however, the window size $W_{ij}$ is selected adaptively based on local homogeneity. Because the final test statistic is computed using a data-dependent window size, its null distribution (and therefore $p$-value calibration) may differ from the fixed-window setting; this point merits clarification and validation.

- Recommendations

Please discuss whether a fixed-window null distribution is an appropriate approximation for the adaptive-window statistic. Ideally, the calibration procedure should incorporate the adaptive selection mechanism (e.g., by simulating the full adaptive procedure under $H_0$ and estimating the null distribution of the final standardized statistic). If a full simulation is too expensive, a practical alternative is to calibrate separately for each admissible window size and to standardize per pixel using its final window size.
\end{reviewbox}

\begin{responsebox}{Response}

\assigned{Revision assigned to Janeth}

\end{responsebox}

\end{document}

